% Miller-Rabin Primality Test \& Pollard's Rho

Deterministic 64-bit primality check in $O(\log N)$ and factorization in $O(N^{1/4})$.

\begin{lstlisting}
typedef unsigned long long ull;

ull mul_mod(ull a, ull b, ull m) {
    long long res = a * b - m * (ull)((long double)a * b / m);
    return (res < 0 ? res + m : res) % m;
}

ull power(ull base, ull p, ull mod) {
    ull res = 1; base %= mod;
    while (p > 0) {
        if (p & 1) res = mul_mod(res, base, mod);
        base = mul_mod(base, base, mod); p >>= 1;
    }
    return res;
}

bool miller_rabin(ull n) {
    if (n < 2) return false;
    if (n == 2 || n == 3) return true;
    if (n % 2 == 0) return false;
    ull d = n - 1; int s = 0;
    while (d % 2 == 0) { d /= 2; s++; }
    static const ull bases[] = {2, 3, 5, 7, 11, 13, 17, 19, 23, 29, 31, 37};
    for (ull a : bases) {
        if (n <= a) break;
        ull x = power(a, d, n);
        if (x == 1 || x == n - 1) continue;
        bool composite = true;
        for (int r = 1; r < s; r++) {
            x = mul_mod(x, x, n);
            if (x == n - 1) { composite = false; break; }
        }
        if (composite) return false;
    }
    return true;
}

ull pollard_rho(ull n) {
    if (n % 2 == 0) return 2;
    if (miller_rabin(n)) return n;
    ull x = 2, y = 2, d = 1, c = 1;
    auto f = [&](ull x, ull n, ull c) { return (mul_mod(x, x, n) + c) % n; };
    while (d == 1) {
        x = f(x, n, c); y = f(f(y, n, c), n, c);
        d = std::gcd(x > y ? x - y : y - x, n);
        if (d == n) { x = rand() % (n - 2) + 2; y = x; c = rand() % (n - 1) + 1; d = 1; }
    }
    return d;
}

void factorize(ull n, vector<ull>& factors) {
    if (n == 1) return;
    if (miller_rabin(n)) { factors.push_back(n); return; }
    ull p = pollard_rho(n); factorize(p, factors); factorize(n / p, factors);
}
\end{lstlisting}
